\documentclass[11pt]{article}
\usepackage{amsmath,amssymb,amsthm}
\newtheorem{theorem}{Theorem}
\title{A Sixteen-Wire Common-Pool Obstruction}
\author{}\date{}\begin{document}\maketitle
\paragraph{Status and scope.}
This is an exact finite and small-parameter obstruction to a single
common residue vector serving every butterfly target, with separate
block biases for the targets. The residues may have either sign.
It does not obstruct target-adaptive residues, other pole or query
geometries, or the original HeadComplexity model. Verification is
written plus exact rational and Laurent arithmetic, not Lean.

Use the notation and query convention of
\texttt{butterfly-adaptive-haar-biased.tex}, with sixteen wires and
\[
 p_y(\epsilon)=y_0\epsilon+y_1\epsilon^2+y_2\epsilon^4+y_3\epsilon^8,
 \qquad t_{k,s,g}=1-p_s.
\]
\begin{theorem}
For every sufficiently small positive $\epsilon$, no one vector
$c\in\mathbb R^{16}$ satisfies every butterfly target using separate
block biases. The same statement holds at $\epsilon=1/100$.
\end{theorem}
\begin{proof}
Eliminating a block bias gives a strict comparison between each
positive-labelled prefix $a$ and negative-labelled prefix $b$.
After dividing by $|p_a-p_b|$, this is $r\cdot c>0$, where
\[
 r_y=\frac{u_y}{(1-p_a-p_y)(1-p_b-p_y)},\qquad
 u_y=\operatorname{sgn}(p_a-p_b)
       u_{k,g}(\pi^{-1}(y)).                                  \tag{1}
\]
The file \texttt{butterfly-common-pool-h16-dual.json} gives seventeen
such rows, possibly coming from different target words, and strictly
positive rational weights at $\epsilon=1/100$ summing to one whose
weighted row sum is exactly zero. Multiplying seventeen required
strict inequalities by these weights gives a contradiction.

The independent checker
\texttt{butterfly-common-pool-h16-dual-check.py} first reconstructs a
butterfly target for every saved row and verifies its actual permutation.
Here is the construction behind that check. A row at level $k$ has one
nonzero entry in each of the $2^{k+1}$ low-bit classes. Its signs determine
the level-$k$ target bits. Earlier layers are set to the identity.
The selected representative in each low-bit class determines all desired
higher output bits. Later switches route each representative to those
higher bits independently, because distinct low-bit classes never reach
the same later switch. The checker also verifies that $a$ is labelled
positive and $b$ negative. Thus all seventeen inequalities are necessary
for a common pool; no infeasibility of an enlarged artificial system is
being substituted for this requirement.

For the small-parameter statement, solve the seventeen equations
\[
 \sum_{i=1}^{17}\lambda_i r_i=0,\qquad
 \sum_{i=1}^{17}\lambda_i=1                                  \tag{2}
\]
over $\mathbb Q(\epsilon)$. Exact guarded power-series elimination,
with tracked precision $1024$, has fourteen pivots of valuation zero
and three of valuation nine. It gives the following leading terms,
in the row order stored in the certificate:
\[
\begin{gathered}
 \tfrac15\epsilon^{10},\ \tfrac3{10}\epsilon^{10},\
 \tfrac2{25}\epsilon^2,\ \tfrac3{10}\epsilon,\
 \tfrac3{10}\epsilon,\ \tfrac12,\ \tfrac12\epsilon,\
 \tfrac15\epsilon,\ \tfrac3{10}\epsilon,\
 \tfrac3{10}\epsilon,\ \tfrac12,\ \tfrac12\epsilon^{10},\
 \tfrac12\epsilon^{10},\ \tfrac1{10}\epsilon^{10},\
 \tfrac3{10}\epsilon,\ \tfrac2{25}\epsilon^2,\
 \tfrac3{10}\epsilon.
\end{gathered}
\]
Every leading coefficient is positive. Nonzero pivots certify an
invertible rational-function system; the checker substitutes its
solution into every equation in (2). Tracked precision certifies the
listed nonzero leading terms. Therefore all seventeen rational weights
are strictly positive on a common punctured interval to the right of
zero. The contradiction from (1) applies throughout that interval.
\end{proof}

\paragraph{Consequence for the proposed induction.}
A common strictly positive eight-wire pool does exist, as certified by
\texttt{butterfly-adaptive-haar-common-positive-h8-laurent.json}.
Its openness gives the written universal sixteen-wire decoder with
residues depending on the target. The theorem above explains why
repeating that argument with a strict common sixteen-wire pool fails
on this particular pole grid. It does not exclude weaker pairwise
cone bridges, additional query profiles, or singular perturbations
with target-dependent residues.
\end{document}
